% BSDE Connection: CX Dealer Market as a Jump-BSDE
% This section establishes that the Cont-Xiong dealer market model
% is equivalent to a backward stochastic differential equation with
% Poisson jumps (BSDEJ), connecting the numerical implementation to
% the BSDE theory of Pardoux-Peng (1990) and El Karoui et al. (1997).

\section{The Dealer Market as a BSDEJ}
\label{sec:bsde-connection}

We show that the Cont--Xiong dealer market model, formulated as a
system of coupled Hamilton--Jacobi--Bellman equations in their
Proposition~3.4, is equivalent to the solution of a backward stochastic
differential equation with Poisson jumps (BSDEJ). This connection places
the numerical methods developed in Chapters~4--5 within the theoretical
framework of BSDE theory and justifies the ``deep BSDE'' methodology.

\subsection{Forward-Backward System}

Consider a single representative market maker (in the mean-field limit)
with inventory process $q_t$ evolving as a compound Poisson process:
\begin{equation}\label{eq:forward-sde}
  dq_t = \Delta \bigl(dN^b_t - dN^a_t\bigr),
\end{equation}
where $N^a_t$ and $N^b_t$ are independent Poisson counting processes
with stochastic intensities
\begin{equation}
  \nu^a_t = \lambda^a f^a\bigl(\delta^a_t, \mu_t\bigr), \qquad
  \nu^b_t = \lambda^b f^b\bigl(\delta^b_t, \mu_t\bigr),
\end{equation}
and $f^a, f^b$ are the Cont--Xiong execution probabilities (their
equation~6), which depend on the agent's own quote $\delta$ and the
population distribution of quotes $\mu_t$.

The market maker's objective over a finite horizon $[0, T]$ is:
\begin{equation}\label{eq:objective}
  V(t, q) = \sup_{\delta^a, \delta^b} \mathbb{E}\Bigl[
    \int_t^T e^{-r(s-t)} \bigl(\text{profits}_s - \psi(q_s)\bigr)\,ds
    + e^{-r(T-t)} g(q_T)
  \Bigr],
\end{equation}
where $\psi(q) = \phi q^2$ is the running inventory penalty and
$g(q_T) = -\psi(q_T)$ is the terminal condition.

\subsection{BSDEJ Formulation}

By the nonlinear Feynman--Kac theorem for jump-diffusion processes
\cite{barles1997}, the value function $V(t, q)$ is characterised as
the $Y$-process of the BSDEJ:
\begin{equation}\label{eq:bsdej}
  -dY_t = f(t, q_t, Y_t, U^a_t, U^b_t)\,dt
          - U^a_t \bigl(dN^a_t - \nu^a_t\,dt\bigr)
          - U^b_t \bigl(dN^b_t - \nu^b_t\,dt\bigr),
\end{equation}
with terminal condition $Y_T = g(q_T) = -\psi(q_T)$, where:
\begin{itemize}
  \item $Y_t = V(t, q_t)$ is the value process,
  \item $U^a_t = V(t, q_t - \Delta) - V(t, q_t)$ is the jump coefficient
        for ask executions (value change when inventory decreases),
  \item $U^b_t = V(t, q_t + \Delta) - V(t, q_t)$ is the jump coefficient
        for bid executions (value change when inventory increases).
\end{itemize}

The generator $f$ encodes the optimal market-making strategy:
\begin{equation}\label{eq:generator}
  f(t, q, Y, U^a, U^b) = -rY - \psi(q)
    + \lambda^a \sup_{\delta^a} \bigl[f^a(\delta^a, \mu)\,(\delta^a \Delta + U^a)\bigr]
    + \lambda^b \sup_{\delta^b} \bigl[f^b(\delta^b, \mu)\,(\delta^b \Delta + U^b)\bigr].
\end{equation}

\begin{remark}[Connection to the HJB system]
The Cont--Xiong Bellman equation (their equation~28) is the
\emph{stationary} ($T \to \infty$) limit of~\eqref{eq:bsdej}. In the
stationary case, $\partial_t V = 0$ and $Y_t = V(q_t)$ depends only on
inventory, not time. The BSDEJ reduces to the algebraic system
\begin{equation}
  rV(q) + \psi(q) = \lambda^a \sup_{\delta^a}\bigl[f^a\,(\delta^a \Delta + U^a(q))\bigr]
                   + \lambda^b \sup_{\delta^b}\bigl[f^b\,(\delta^b \Delta + U^b(q))\bigr],
\end{equation}
which is exactly the system solved by Algorithm~1 and our neural
Bellman solver. The neural solver learns $V(q)$; the jump coefficients
$U^a(q) = V(q-1) - V(q)$ and $U^b(q) = V(q+1) - V(q)$ are computed
from value differences; and optimal quotes are derived from the
first-order condition of~\eqref{eq:generator}.
\end{remark}

\begin{remark}[Absence of diffusion term]
Unlike the standard BSDE framework where $Z_t\,dW_t$ represents the
diffusion coefficient, the CX model has $Z_t = 0$ because the value
function does not depend on the mid-price $S_t$ (the price cancels
in the Cont--Xiong formulation; see their Lemma~2.6). All stochasticity
enters through the Poisson jump terms $U^a\,dN^a + U^b\,dN^b$. This
makes the problem a \emph{pure-jump BSDE}, which is well-posed under
weaker conditions than the general BSDE theory requires
\cite{tang1994}.
\end{remark}

\subsection{Deep BSDE Method for Jump Processes}

The deep BSDE method of \citet{han2018} parameterises $Z_t$ as neural
networks at each time step and propagates the BSDE forward in time,
minimising the terminal mismatch $|Y_T - g(q_T)|^2$. For our
pure-jump problem, we parameterise $U^a_t$ and $U^b_t$ instead of
$Z_t$:

\begin{enumerate}
  \item \textbf{Initialise:} Learn $Y_0 = V(0, q_0)$ as a scalar parameter.
  \item \textbf{Forward propagation:} At each time step $t_m$, the subnet
        $\phi_m(t, q) \to (U^a_m, U^b_m)$ predicts the jump coefficients.
        Optimal quotes $\delta^{a*}_m, \delta^{b*}_m$ are derived from the
        FOC of~\eqref{eq:generator}.
  \item \textbf{BSDE step:} Update $Y_{m+1} = Y_m - f_m \Delta t +
        U^a_m (N^a_{m+1} - N^a_m - \nu^a_m \Delta t) +
        U^b_m (N^b_{m+1} - N^b_m - \nu^b_m \Delta t)$.
  \item \textbf{Loss:} $\mathcal{L} = \mathbb{E}|Y_M - g(q_T)|^2$.
\end{enumerate}

This is a direct extension of \citet{han2018} to jump processes,
following the theoretical framework of \citet{han2022} for
McKean--Vlasov FBSDEs. The mean-field coupling enters through the
execution probabilities $f^a, f^b$, which depend on the population
distribution $\mu_t$.

\subsection{Relationship to the Stationary Solver}

Our primary numerical results (Chapters~4--5) use the stationary
($T = \infty$) formulation, which we call the ``neural Bellman solver.''
This corresponds to solving the BSDEJ in the ergodic limit:
\begin{equation}
  \lim_{T \to \infty} \frac{1}{T} V(0, q_0; T) = \text{ergodic constant},
\end{equation}
and the optimal quotes converge to the stationary Nash equilibrium.

The finite-horizon BSDEJ solver (this section) serves as validation:
for large $T$, the finite-horizon and stationary solutions should agree
at $t = 0$. We verify this computationally by comparing the deep BSDEJ
output $Y_0$ against the stationary Algorithm~1 value $V(0)$.

\begin{remark}[Why BSDE and not just HJB?]
The BSDE formulation is preferred over the direct HJB approach for
three reasons: (i)~it extends naturally to high-dimensional problems
via the deep BSDE method, avoiding the curse of dimensionality that
afflicts grid-based HJB solvers; (ii)~the Wasserstein convergence
guarantees of \citet{han2022} apply to the BSDE formulation but not
directly to the HJB; (iii)~the probabilistic interpretation enables
simulation-based training without discretising the state space.
\end{remark}
